\pagestyle{plain}


\chapter{Referencial teórico}


\section{Desenvolvimento Rural e Sustentável: Perspectivas e Desafios no Brasil}


O conceito de desenvolvimento vai além do crescimento econômico, englobando aspectos socioculturais como saúde, educação, qualidade de vida e bem-esta (Sen, 1999). No Brasil, essa discussão ganha contornos específicos quando voltada ao desenvolvimento rural, um fenômeno que engloba não apenas crescimento econômico, mas também a melhoria da qualidade de vida, as relações sociais e a autonomia econômica das comunidades rurais \cite{Navarro2001}. 

Segundo \citeonline{Abramovay2020}, o desenvolvimento rural brasileiro segundo a teoria seniana, desenvolvimento como liberdade, ainda carece de um horizonte estratégico claro que integre ações públicas e privadas com vistas a ampliar as liberdades substantivas dos indivíduos enquanto preserva e regenera os ecossistemas dos quais a sociedade depende. Esse desenvolvimento rural sustentável envolve além das intitulações financeiras  a diversificação dos meios de vida, onde o foco não é apenas melhorar a produtividade agrícola, mas fortalecer as capacidades e resiliência das comunidades rurais frente a adversidades econômicas e ambientais \cite{Perondi2012}. 

Assim, em vez de as políticas desenvolvimentistas para o meio rural focarem apenas indicadores econômicos tradicionais, como a renda, o desenvolvimento deve contemplar também aspectos como saúde, educação, acesso à justiça e condições de vida adequadas. Segundo \citeonline{Sen1999}, o desenvolvimento autêntico ocorre quando as pessoas possuem as liberdades substantivas para escolher e conduzir suas vidas de maneira autônoma, uma perspectiva que influenciou políticas públicas e iniciativas de desenvolvimento em diversas partes do mundo. Assim, alcançar um desenvolvimento rural sustentável e duradouro requer políticas adaptadas às realidades locais, que integrem os diversos meios de subsistência e valorizem os conhecimentos específicos de cada comunidade.


\section{Papel dos Povos Originários e Comunidades Tradicionais na Sustentabilidade e Conservação Ambiental}

Neste contexto, a integração dos Povos originários e Comunidades Tradicionais no processo de desenvolvimento rural é crucial \cite{Abramovay2020}. Estas comunidades são agrupamentos sociais que mantêm uma conexão intrínseca com seus territórios, os quais são transformados em espaços coletivos significativos \cite{Kovaleva2021}. A interação desses grupos com o ambiente e com a sociedade em geral enriquece o seu saber tradicional, contribuindo para o desenvolvimento rural sustentável \cite{Montanari2020}.

O desenvolvimento rural sustentável é fundamental nesse contexto, atuando como um paradigma que alinha o crescimento econômico à preservação ambiental e à justiça social. \citeonline{Sachs2002} amplia essa visão, argumentando que o desenvolvimento rural sustentável pode ser alcançado por meio da exploração responsável dos "três Bs": biomassa, biodiversidade e biotecnologia sustentável. Esses elementos são particularmente relevantes nas comunidades rurais e tradicionais, cuja sobrevivência e crescimento dependem da capacidade de gerir de forma sustentável os recursos naturais \cite{Tagliapietra2021}.

A abordagem centrada nos três Bs permite a construção de uma civilização moderna do vegetal em países tropicais, como o Brasil, onde a interação harmoniosa entre práticas agrícolas e respeito pelo meio ambiente é essencial para o bem-estar e a resiliência das comunidades rurais \cite{Sachs1993}.

Os Povos originários e Comunidades Tradicionais são agrupamentos sociais localizados que estabelecem uma conexão profunda e duradoura com um espaço físico específico, transformando-o em um território coletivo \cite{Grant2020}. Essa transformação é fruto do trabalho dos fundadores das comunidades, que estabeleceram suas raízes nesse local. As comunidades possuem um corpo de conhecimentos e práticas singulares, derivados das interações complexas com o ambiente natural \cite{Clayton2019}. Esse saber é enriquecido tanto pela tradição quanto pela interação com dinâmicas sociais mais amplas com outros grupos sociais e com o meio ambiente.

Os conhecimentos etnoecológicos e etnopedológicos e etnoclimatológico, assim como os saberes agroecológicos das comunidades tradicionais, são aspectos fundamentais de sua relação com o meio ambiente, constituindo pilares da sua identidade cultural, ancestralidade e a sustentabilidade. Esses conhecimentos, como destacado por \citeonline{Motti2023}, englobam a compreensão detalhada dos sistemas ecológicos locais, baseada em séculos de observação e interação direta com o ambiente natural. Esta sabedoria inclui, mas não se limita a técnicas agroecológicas adaptadas, práticas de manejo de recursos naturais, identificação e uso de plantas medicinais, assim como padrões climáticos e comportamento animal \cite{Ahad2023}.

Este entendimento intrincado da natureza, que vai além do conhecimento técnico, é incorporado nas crenças, rituais e narrativas culturais das comunidades, evidenciando uma visão de mundo que harmoniza o uso e a conservação dos recursos naturais. Tais saberes, portanto, não apenas sustentam a subsistência física das comunidades tradicionais, mas também reforçam sua coesão social, resiliência e resistência diante das mudanças ambientais e sociais \cite{Brzi2023}.

Com efeito, a cultura imaterial do rural e ambiental das CPT, entrelaçada com os espaços que ocupam física e simbolicamente, é um testamento vivo de seus direitos sociais \cite{BobatoStadler2020}. Este reconhecimento de seus direitos e de seu espaço tanto no sentido físico quanto no cultural, marca um avanço significativo na valorização e proteção dessas comunidades \cite{Colao2011}.
A autonomia das comunidades tradicionais é fundamental para a manutenção da sustentabilidade de seus membros e preservação da identidade social. Essa autonomia é exercida pela troca cíclica dos Saberes Agroecológicos Tradicionais (SATs) que contribuem para a subsistência, melhorando a sua qualidade de vida, fornecendo alimentos seguros e aumentando a sua resiliência aos recursos naturais \cite{MohdSalim2023}. 

Vale destacar que na década de 1980, no contexto brasileiro, as populações residentes em Áreas Protegidas (APs) frequentemente eram subvalorizadas ou ignoradas nas políticas ambientais \cite{Marques2016}. De acordo com \citeonline{Thum2017}, desde a última década do século XX, a questão das populações tradicionais já se fazia presente nos debates sobre a diferença de modos de ser de populações. No decorrer desse processo, diferentes tentativas de definição foram colocadas em circulação. Naquele momento se formulava uma definição com caracterização capaz de incorporar as diferentes dimensões presentes nesses segmentos.

Embora, historicamente chamados planeadores ambientais considerem as comunidades locais como a principal ameaça à conservação da biodiversidade ou um retrocesso no desenvolvimento nacional \cite{Singh2000}, tal percepção negligencia o papel integral que tais comunidades desempenham na conservação da biodiversidade. Essas percepções começam a se transformar através de uma sequência de debates, embates e movimentos que culminaram no reconhecimento progressivo da importância conservacionista exercida pelas comunidades, as quais passaram a ser reconhecidas como protetoras da floresta \cite{Asante2017}.


\section{Legislação, Desafios de Implementação e Salvaguarda dos Saberes Tradicionais}


No ano de 2007, o processo de reconhecimento dos povos e comunidades tradicionais (PCTs) no Brasil atingiu um importante marco com a promulgação do Decreto nº 6.040/07 (Institui a Política Nacional de Desenvolvimento Sustentável dos Povos e Comunidades Tradicionais., 2007). Ao estabelecer um quadro legal, o decreto garante e assegura os direitos dessas comunidades ao acesso à terra, marcando um passo fundamental na consolidação de sua presença e relevância no âmbito nacional \cite{Yustiningrum2023}.

Notavelmente, o Artigo 3º, Inciso XV que dispõe ser objetivo específico da PNPCT XV – “reconhecer, proteger e promover os direitos dos povos e comunidades tradicionais sobre os seus conhecimentos, práticas e usos tradicionais” (Institui a Política Nacional de Desenvolvimento Sustentável dos Povos e Comunidades Tradicionais., 2007). Essa disposição legal é um reconhecimento de que a integração de práticas tradicionais e conhecimentos locais na gestão ambiental oferece uma perspectiva mais holística e eficaz para a conservação ambiental \cite{Ciocco2023}. Ao reconhecer formalmente estes conhecimentos e práticas, a legislação proporciona uma base para a sua proteção e promoção, assegurando que eles continuem a contribuir para a sustentabilidade ambiental e cultural \cite{Reyes-Garca2019}. 

A Lei n.º 9.985/2000, que institui o Sistema Nacional de Unidades de Conservação da Natureza (SNUC), reforça essa proteção ao estabelecer diretrizes para o uso sustentável e a preservação dos ecossistemas naturais, incluindo o reconhecimento do valor dos conhecimentos tradicionais para a conservação  \cite{Brasil2015}. Entretanto, a promulgação da Lei n.º 13.123/2015 revogou a Medida Provisória n.º 2.186-16/2001, deixando lacunas em relação à criação de um sistema formal de registro desses conhecimentos tradicionais, particularmente no que se refere ao conhecimento agrícola e outros saberes sui generis \cite{Gomes2016}.

Embora a legislação tenha avançado fortemente, ela não especifica claramente como implementar essas políticas, revelando uma lacuna na sua operacionalização \cite{Tsai2020}. Isso aponta para a necessidade de desenvolver estratégias e ações concretas que integrem e valorizem os conhecimentos e práticas das comunidades tradicionais na gestão ambiental e conservação da biodiversidade, reconhecendo sua autoria e originalidade.

Neste sentido, a proteção dos conhecimentos e saberes agroecológicos tradicionais por mecanismos de registro de conhecimento intelectual e a sua consequente transformação em reconhecimento de originalidade, além de poder solucionar distorções no que diz respeito à importância de tais expressões no plano simbólico para as comunidades tradicionais, introduzem nas mesmas uma lógica de autoria coletiva de determinada comunidade tradicional \cite{LimaVerdanRangel2018}.

Atualmente, conforme dados do Instituto Brasileiro de Geografia e Estatística \cite{IBGE2010}, o Território de identidade semiárido nordeste II abriga 18 municípios com a população total = 400 mil habitantes, estes municípios abrigam cerca de 12 comunidades quilombolas identificadas por registros administrativos e 9 agrupamentos quilombolas \cite{IBGE2022}. Essas comunidades desempenham um papel significativo em nossa memória histórica e cultural do estado (Santos; Santos, 2020). Enfatizar a relevância da Bahia nesse cenário é fundamental para promover a conscientização sobre a contribuição significativa dessas comunidades para a riqueza cultural e agroecológica do estado e do país como um todo \cite{Lima2010}.


\section{Valoração de Ativos Intangíveis e Propriedade Intelectual}

\subsection{Frameworks de Valoração de Conhecimentos Tradicionais}

A valoração de conhecimentos tradicionais como ativos intangíveis representa desafio metodológico que demanda frameworks adaptados às especificidades desses saberes. A fundamentação teórica para compreensão de conhecimento como ativo estratégico remonta à \textbf{Resource-Based View} \cite{Penrose1959,Barney1991}, que posiciona recursos intangíveis, especialmente conhecimento organizacional, como fonte de vantagem competitiva sustentável. A teoria evolucionária de Nelson e Winter \cite{Nelson1982} complementa essa perspectiva ao conceituar conhecimento como rotinas organizacionais incorporadas (embedded), path-dependent e tacitamente transmitidas, características que se aplicam diretamente a saberes tradicionais comunitários.

A literatura sobre gestão do conhecimento \cite{Nonaka1995,Grant1996} distingue conhecimento tácito (implícito, contextual, difícil de codificar) de conhecimento explícito (codificado, transferível, documentado). A gestão estratégica de conhecimentos tradicionais demanda processos de externalização que convertam saberes tácitos em formas explícitas sem descaracterizar sua essência cultural, viabilizando apropriação de valor pelas comunidades detentoras. A teoria da criação de conhecimento de Nonaka e Takeuchi \cite{Nonaka1995}, baseada na espiral SECI (socialização, externalização, combinação, internalização), fornece arcabouço conceitual para essa conversão.

A valoração de ativos intangíveis evoluiu de abordagens tradicionais unidimensionais \cite{SmithParr2000} para frameworks analíticos sofisticados que incorporam múltiplas variáveis e técnicas computacionais avançadas. Pioneiros em capital intelectual como Sveiby \cite{Sveiby1997}, Stewart \cite{Stewart1997} e Edvinsson \cite{Edvinsson1997} estabeleceram taxonomias que classificam ativos intangíveis em capital humano (competências individuais), capital estrutural (processos, sistemas, propriedade intelectual) e capital relacional (redes, parcerias, reputação). Essas categorias, quando adaptadas ao contexto comunitário, permitem mensuração sistemática de diferentes dimensões do valor de conhecimentos tradicionais.

Smith e Parr \cite{SmithParr2000} consolidaram três abordagens clássicas de valoração. O método do custo estima valor baseado no investimento acumulado necessário para recriar o ativo, enquanto o método de mercado utiliza comparação com transações similares de ativos correlatos. O método da renda, por sua vez, projeta fluxos de caixa futuros descontados pela exploração comercial do ativo. O Balanced Scorecard \cite{Kaplan1996} complementa essas abordagens ao integrar perspectivas financeiras, de clientes, de processos internos e de aprendizado/crescimento, viabilizando mensuração holística de valor estratégico.

Zhou e Wang (2022) propõem avanço metodológico ao integrar métodos paramétricos com algoritmos de machine learning para valoração de ativos intangíveis. A abordagem paramétrica utiliza modelos estatísticos baseados em variáveis-chave observáveis, enquanto machine learning identifica padrões complexos e interações não-lineares que métodos tradicionais não capturam. Algoritmos como random forests e redes neurais superam modelos paramétricos na previsão de valor ao processar simultaneamente características tecnológicas, de mercado e organizacionais.

O modelo VAIC (Value Added Intellectual Coefficient) \cite{Pulic2000} oferece ferramenta analítica que mensura a eficiência com que capital intelectual é convertido em valor econômico. O modelo é expresso pela formulação VAIC = HCE + SCE + CEE, em que HCE (Human Capital Efficiency) representa a eficiência do capital humano na geração de valor adicionado; SCE (Structural Capital Efficiency) expressa a eficiência das estruturas e processos organizacionais; CEE (Capital Employed Efficiency) mede a eficiência do capital físico e financeiro. Adaptado ao contexto de comunidades tradicionais, esse framework permite mensurar como conhecimentos, práticas organizativas comunitárias e recursos materiais convertem-se em valor econômico e social.

\subsection{Mercados para Tecnologia e Estratégias de Comercialização}

A comercialização de conhecimentos tradicionais como ativos tecnológicos insere-se na dinâmica dos mercados para tecnologia \cite{AroraEtAl2001}, caracterizados por divisão do trabalho inovativo entre entidades especializadas em gerar conhecimento versus aquelas focadas em comercialização. Arora, Fosfuri e Gambardella (2001) demonstram que a eficiência desses mercados depende criticamente da estrutura de apropriabilidade e dos mecanismos de redução de assimetrias informacionais. Em mercados espessos (thick markets), onde múltiplas transações ocorrem com frequência, métodos comparativos de mercado são viáveis; em mercados finos (thin markets), típicos de conhecimentos especializados, métodos baseados em renda descontada tornam-se preferíveis.

Wu e Li (2022) complementam essa perspectiva ao desenvolver modelo de contribuição tecnológica que distingue contribuições quantitativas (métricas objetivas de performance) de contribuições qualitativas (valor estratégico, posicionamento competitivo e análise de risco). No contexto de conhecimentos tradicionais agroecológicos, contribuições quantitativas incluem eficiência no uso de recursos, produtividade agrícola e resiliência climática, enquanto contribuições qualitativas abrangem valor cultural, potencial de diferenciação mercadológica e sustentabilidade socioambiental.

\section{Machine Learning na Salvaguarda dos Saberes Tradicionais Agroecológicos}

O aprendizado de máquina (ML), especialmente em contextos agroecológicos, desempenha um papel crítico ao modelar interações complexas entre variáveis ambientais e socioeconômicas \cite{Zhou2018}. Em ML supervisionado, algoritmos são treinados em dados rotulados, aprendendo a prever respostas específicas (outputs) a partir de um conjunto de variáveis de entrada (inputs) conhecido. No contexto de gestão de conhecimentos tradicionais, ML oferece ferramentas que abrangem desde a identificação automatizada de padrões em grandes volumes de dados etnográficos até a classificação de práticas agroecológicas segundo taxonomias validadas, passando pela predição de eficácia de saberes tradicionais em diferentes contextos ambientais e pela valoração de ativos intangíveis baseada em análise de múltiplas variáveis simultâneas. 

Para isso, a seleção de variáveis é otimizada por técnicas de redução dimensional, que eliminam redundâncias e priorizam variáveis de maior relevância para o fenômeno estudado, minimizando o risco de sobreajuste (overfitting) e mantendo a interpretabilidade. Métodos como florestas aleatórias e máquinas de vetor de suporte não apenas identificam variáveis críticas, mas também hierarquizam sua importância relativa, fornecendo insights sobre a influência de cada variável no contexto agroecológico.

Neste sentido, o uso de machine learning permite a realização de treinamentos de modelos a partir de dados empíricos, aprimorando a capacidade de prever e identificar padrões em novos conjuntos de dados  \cite{Choudhury2021}, questões estas que por vezes não são alcançadas por análises estatísticas de hipóteses. 
Em contraste aos métodos de hipóteses, os métodos de aprendizado de máquina (ML) podem ser aplicados em pesquisas guiadas pela descoberta, como abordagens indutivas ou abdutivas \cite{Ritchie2016}. Isso ocorre porque, ao contrário dos métodos tradicionais, os algoritmos de ML são capazes de identificar padrões complexos em variáveis explicativas ($X$) que se relacionam com o resultado ($y$), explorando estruturas que não foram previamente especificadas \cite{Benos2021}. 

Diferentemente dos testes de hipótese, os algoritmos de ML constroem modelos com formas funcionais flexíveis, otimizando o desempenho ao usar as variáveis explicativas ($X$) para prever o resultado ($y$) \cite{Parish2016}. Essas formas funcionais resultantes podem evidenciar relações subjacentes e inesperadas nos dados.
Em outras palavras, em vez de testar dedutivamente um modelo previamente estabelecido pelo pesquisador, como ocorre nas análises comportamentais tradicionais focada em inferência, os algoritmos de ML constroem um modelo indutivamente, com base nos próprios dados, para revelar padrões latentes \cite{Candelieri2019}. Essas características tornam o ML também uma ferramenta eficaz para a análise pós-hoc de resultados de regressão tradicionais, permitindo a identificação de padrões que poderiam passar despercebidos em uma abordagem convencional \cite{vanVliet2020}.

\section{Gestão da Inovação Aplicada a Conhecimentos Tradicionais}

\subsection{Fundamentos da Gestão da Inovação}

A gestão da inovação emerge como disciplina estratégica fundamental no contexto contemporâneo, caracterizado pela aceleração tecnológica e pela necessidade de adaptação contínua das organizações \cite{Tidd2005}. Conforme definido por Tidd, Bessant e Pavitt, a gestão da inovação envolve processos sistematizados de busca, seleção, implementação e captura de valor de oportunidades inovadoras, operacionalizados através de quatro fases críticas, iniciando pelo \textit{scanning} (detecção de oportunidades e ameaças no ambiente), seguido pela \textit{strategy formulation} (definição de opções estratégicas), pela \textit{resource allocation} (alocação de recursos) e pela \textit{implementation and learning} (implementação e aprendizado iterativo). 

Este framework teórico, originalmente desenvolvido para contextos organizacionais convencionais, apresenta aplicabilidade significativa ao domínio dos conhecimentos tradicionais quando adaptado para realidades de comunidades e sistemas organizacionais descentralizados. A transposição conceitual demanda reconhecimento explícito de que conhecimentos tradicionais não constituem estruturas estáticas, mas sim sistemas adaptativos complexos que evoluem através de mecanismos endógenos de inovação e aprendizado comunitário.

\subsection{Paradigmas de Inovação e Conhecimentos Tradicionais}

\subsubsection{Inovação Aberta e Integração de Saberes}

O paradigma da \textbf{inovação aberta} (\textit{open innovation}), proposição conceitual de Chesbrough \cite{Chesbrough2003}, estrutura-se em três processos interdependentes, compreendendo o \textit{outside-in process} (integração de conhecimentos externos), o \textit{inside-out process} (circulação de conhecimentos internos) e o \textit{coupled process} (co-criação com parceiros externos). Este framework demonstra particular relevância para contextos de conhecimentos tradicionais, onde a integração simbiótica entre saberes endógenos comunitários e conhecimentos científicos contemporâneos potencializa oportunidades inovadoras que transcendem as capacidades isoladas de cada fonte de conhecimento.

No contexto de comunidades tradicionais, o outside-in process materializa-se através da absorção seletiva de tecnologias digitais, metodologias de pesquisa e ferramentas de gestão de propriedade intelectual, enquanto o inside-out process caracteriza-se pela documentação, sistematização e oferta de saberes tradicionais para parcerias colaborativas com instituições de pesquisa. O coupled process, particularmente relevante para inovação social, opera mediante co-criação entre comunidades e pesquisadores, viabilizando conversão de conhecimentos tradicionais em tecnologias sociais inovadoras que preservam autenticidade cultural enquanto agregam aplicabilidade científica.

\subsubsection{Inovação de Usuário e Agência Comunitária}

Von Hippel \cite{vonHippel1988} estabeleceu arcabouço conceitual fundamental ao demonstrar que inovação frequentemente origina-se de usuários que detêm conhecimento contextual especializado, acesso a informações de mercado e motivação intrínseca para resolver problemas específicos de seu domínio. Esta proposição aplica-se com particular pertinência a comunidades tradicionais, que constituem expertise holders primários em sistemas de manejo de recursos naturais, práticas agroecológicas e conhecimentos etnoecológicos. A reconceptualização de comunidades tradicionais como agentes inovadores, e não meramente como detentores passivos de patrimônio cultural, reposiciona dinamicamente o locus da inovação, deslocando-o de instituições formais de pesquisa para territórios de prática comunitária.

\subsection{Gestão Estratégica de Conhecimentos Tradicionais como Ativos Intangíveis}

\subsubsection{Resource-Based View e Propriedade Intelectual Coletiva}

A \textbf{Resource-Based View} \cite{Barney1991}, arcabouço analítico para compreensão de vantagem competitiva sustentável, postula que recursos valiosos, raros, inimitáveis e não-substituíveis (VRIN) constituem fundamento de diferencial competitivo duradouro. Conhecimentos tradicionais, quando adequadamente protegidos e geridos, atendem integralmente aos critérios VRIN. São \textbf{valiosos} pelo potencial de aplicação em inovação agrícola, farmacológica e de manejo ambiental, e \textbf{raros} por resultarem de séculos de observação sistemática e experimentação contextualizada. Caracterizam-se ainda como \textbf{inimitáveis}, por originarem-se de processos de aprendizado \textit{path-dependent} acumulados em contextos ecológicos e culturais específicos, e \textbf{não-substituíveis}, por constituírem fonte primária e irreplicável de conhecimentos sobre biodiversidade e sustentabilidade territorial.

A apropriação de valor gerado por conhecimentos tradicionais demanda, contudo, frameworks de gestão que transcendam modelos convencionais de propriedade intelectual individual. A \textbf{propriedade intelectual coletiva}, materializada através de indicações geográficas, marcas coletivas e proteções sui generis, constitui instrumento jurídico e estratégico essencial para garantir que benefícios econômicos derivados da exploração comercial de conhecimentos tradicionais revertem-se às comunidades detentoras, assegurando apropriação justa e sustentabilidade de iniciativas de inovação baseadas em saberes endógenos.

\subsubsection{Capacidade Absortiva Comunitária}

Cohen e Levinthal \cite{CohenLevinthal1990} desenvolveram conceito fundamental de \textbf{capacidade absortiva}, definida como habilidade organizacional de reconhecer valor de informação externa, assimilá-la e aplicá-la para propósitos comerciais. Zahra e George \cite{Zahra2002} refinaram esta concepção ao distinguir capacidade absortiva potencial (aquisição e assimilação de conhecimento externo) de capacidade absortiva realizada (transformação e exploração de conhecimento).

No contexto de comunidades tradicionais, o fortalecimento de capacidade absortiva operacionaliza-se através de múltiplas dimensões. A \textbf{aquisição} envolve o acesso a tecnologias de machine learning, ferramentas de documentação digital e marcos regulatórios de propriedade intelectual. A \textbf{assimilação} refere-se à integração de conhecimentos tecnológicos sem comprometimento de autenticidade cultural. A \textbf{transformação} diz respeito à conversão de saberes tradicionais em representações estruturadas (bancos de dados, modelos ontológicos) que facilitem análise computacional e aplicação. A \textbf{exploração}, por fim, traduz conhecimentos transformados em inovações comercializáveis e escaláveis.

O desenvolvimento cumulativo e path-dependent de capacidade absortiva constitui pré-requisito essencial para que comunidades tradicionais e organizações parceiras possam apropriar-se efetivamente de conhecimentos e tecnologias externas, convertendo-os em inovação social e desenvolvimento territorial sustentável.

\subsection{Ecossistemas de Inovação e Sistemas Agroecológicos Tradicionais}

\subsubsection{Ecossistemas de Inovação Territorial}

Adner \cite{Adner2006} conceitua ecossistema de inovação como estrutura de interdependências tecnológicas entre componentes, onde a viabilidade de qualquer componente depende da disponibilidade de outros. Esta perspectiva revela-se particularmente frutífera para compreensão de sistemas agroecológicos tradicionais em contextos de inovação colaborativa.

A análise de assimetrias tecnológicas upstream (componentes essenciais integrados) e downstream (complementos que agregam valor ao produto final) \cite{AdnerKapoor2010} ilumina estratégias para posicionamento de conhecimentos tradicionais em cadeias produtivas sustentáveis. Comunidades que dominam saberes sobre manejo de recursos naturais, práticas agroecológicas e conhecimentos etnoecológicos controlam componentes tecnológicos upstream de relevância crítica para bioeconomia contemporânea, podendo fortalecer significativamente sua posição negociadora em arranjos de comercialização e parcerias colaborativas.

\subsubsection{Inovação Social e Tecnologias Sociais}

Inovação social caracteriza-se pela alta intensidade de engajamento comunitário, exploração de oportunidades derivadas de conhecimento local especializado e navegação em ambientes de participação coletiva. As tecnologias sociais, definidas como produtos, processos ou serviços que incorporam conhecimento local em soluções inovadoras para problemas socioambientais, constituem materialização operacional de inovação social no contexto de sistemas agroecológicos tradicionais.

A conversão de práticas agroecológicas tradicionais em tecnologias sociais escaláveis demanda documentação sistemática de conhecimentos por meio de protocolos participativos, validação técnico-científica que não comprometa a legitimidade epistemológica comunitária, adaptação de práticas para contextos variados mantendo princípios fundamentais, e desenvolvimento de modelos de sustentabilidade econômica que garantam apropriação comunitária de valor.

\subsection{Capacidades Dinâmicas e Inovação Agroecológica}

\subsubsection{Sensing, Seizing e Reconfiguring}

Teece \cite{Teece2007} propôs framework de capacidades dinâmicas estruturado em três dimensões críticas. O \textbf{sensing} corresponde à capacidade de identificar oportunidades e ameaças no ambiente, o \textbf{seizing} traduz-se na mobilização de recursos organizacionais para capturar oportunidades identificadas, e o \textbf{reconfiguring} refere-se à adaptação e reconfiguração de estruturas organizacionais em resposta a mudanças ambientais e sociais.

No contexto de comunidades tradicionais aplicando inovação agroecológica, estas dimensões materializam-se de forma concreta. O \textbf{sensing} compreende o monitoramento de mudanças climáticas, a identificação de mercados emergentes para produtos agroecológicos e o reconhecimento de oportunidades tecnológicas. O \textbf{seizing} envolve a mobilização de recursos comunitários, o estabelecimento de parcerias com instituições de pesquisa e a implementação de tecnologias digitais para gestão. O \textbf{reconfiguring} abrange o ajuste de sistemas produtivos diante de condições ambientais alteradas, a modificação de estruturas organizacionais para apropriação de oportunidades e a reconfiguração de arranjos institucionais para acesso a mercados.

O desenvolvimento destas capacidades dinâmicas constitui fundamento essencial para que comunidades tradicionais transitem de posições defensivas de preservação de conhecimentos para posições ofensivas de apropriação ativa de oportunidades inovadoras.

\subsection{Valoração Integrada Multidimensional de Conhecimentos Tradicionais}

A valoração de conhecimentos tradicionais transcende abordagens unidimensionais baseadas exclusivamente em métricas econômicas convencionais. Frameworks analíticos contemporâneos reconhecem múltiplas dimensões de valor que devem ser integradas em avaliações compreensivas:

A dimensão \textbf{econômica} constitui-se em potencial de retorno financeiro direto, mensurado através de abordagens de fluxo de caixa descontado, análise de valor presente líquido de ativos intelectuais, e modelagem econométrica de contribuição tecnológica em cadeias produtivas. O modelo VAIC (Value Added Intellectual Coefficient) \cite{Pulic2000} oferece instrumento específico para mensuração de eficiência na conversão de capital intelectual em valor agregado.

A dimensão \textbf{científica} refere-se à contribuição potencial de conhecimentos tradicionais para avanços em compreensão de processos ecológicos, adaptação climática, manejo sustentável de recursos e inovação em campos como etnobotânica, etnoecologia e agroecologia.

A dimensão \textbf{cultural} engloba valor intrínseco de conhecimentos para manutenção da identidade comunitária, transmissão intergeracional de saberes, preservação de patrimônio cultural imaterial e reafirmação de soberania epistemológica. Embora frequentemente negligenciada em frameworks de gestão convencionais, a dimensão cultural constitui fundamento essencial de sustentabilidade de iniciativas de inovação baseadas em conhecimentos tradicionais.

A dimensão \textbf{socioambiental} compreende contribuições de conhecimentos tradicionais para conservação da biodiversidade, estabilização de ecossistemas, adaptação às mudanças climáticas e promoção de resiliência territorial.

A síntese de gestão da inovação, propriedade intelectual coletiva, metodologias participativas e tecnologias computacionais emergentes oferece fundamento para desenvolvimento de modelos integrados que operacionalizam, na prática, conceitos de inovação social, desenvolvimento territorial sustentável e soberania epistemológica comunitária.

Assim, o presente projeto de pesquisa propõe, desenvolver um modelo empírico por meio de machine learning para salvaguarda dos saberes agroecológicos dos povos originários e comunidades tradicionais do Território de Identidade Semiárido Nordeste II na Bahia. Espera-se que o modelo se apresente como uma ferramenta que sustente a identidade cultural e ambiental do semiárido baiano, promovendo a salvaguarda dos direitos de propriedade intelectual e usufruto dos saberes e conhecimentos que respeite e reforce a singularidade desses saberes dos povos detentores.


